\worksheettitle
  {Составление чисел из цифр}
  {\modulemark{M05} Версия учителя \textbullet\ маршрут 30–45 минут}

\begin{teacherbox}[Паспорт]
  \textbf{Результат:} ученик строит полный упорядоченный список с учётом
  повторений и ведущего нуля. \textbf{Опора:} M01 и M02.

  Центральная идея — группировка по старшей цифре. Перестановки и формулы не
  вводятся как новая тема; счёт вариантов используется только для проверки.
\end{teacherbox}

\begin{checkpointbox}[Вход]
  В числе $525$ цифра десятков — 2, её значение 20. Числа $25$ и $52$
  различны, потому что одинаковые цифры занимают разные разряды.
\end{checkpointbox}

\section{Две модели}

С повторениями из 2 и 5: $22,25,52,55$. Без повторений: $25,52$.
Проговаривайте не только ответ, но и изменение условия.

Из цифр 2, 5, 7 без повторений: $257,275$; $527,572$; $725,752$.
Полнота видна по трём группам старшей цифры и двум порядкам остатка.

\begin{teacherbox}[Ключевые вопросы]
  «Какую цифру ставим первой? Какие цифры после этого остались? Все ли
  варианты с этой первой цифрой записаны? Перешли ли мы ко всем допустимым
  первым цифрам?»
\end{teacherbox}

\section{Ноль и ведущая позиция}

Из 0, 3, 8: $308,380,803,830$. Ветвь с нулём в сотнях не создаёт
трёхзначных чисел.

Ответы базы:
\begin{enumerate}
  \item $11,16,61,66$;
  \item $16,19,61,69,91,96$;
  \item $409,490,904,940$.
\end{enumerate}

В ошибочном списке потеряны $33$ и $77$: ученик самовольно применил запрет
повторений.

\section{Как дозировать лист}

\begin{resultbox}[Маршруты]
  \begin{itemize}
    \item \textbf{Быстрый:} две модели, ноль, по одному заданию закрепления,
      проверка чужого списка, контроль, M05\(\star\).
    \item \textbf{Основной:} вся база, 3–4 задания закрепления, сравнение
      условий, контроль.
    \item \textbf{С опорой:} модели с раскладкой карточек, M05-R, затем новые
      короткие списки и контроль.
  \end{itemize}
\end{resultbox}

\section{Закрепление: ответы}

Двузначные:
\begin{enumerate}
  \item $44,48,84,88$;
  \item $24,28,42,48,82,84$;
  \item $50,55,57,70,75,77$ (ноль не может стоять первым);
  \item $50,57,70,75$.
\end{enumerate}

Трёхзначные:
\begin{enumerate}
  \item $136,163,316,361,613,631$;
  \item $205,250,502,520$;
  \item $222,229,292,299,922,929,992,999$.
\end{enumerate}

\begin{teacherbox}[Замечание к формулировке]
  В последнем задании сказано «использовать можно обе цифры», а не «обе цифры
  должны встретиться». Поэтому $222$ и $999$ входят в список. Это намеренная
  проверка внимательного чтения условия.
\end{teacherbox}

\section{Сравнение условий}

Для цифр 1, 4, 7 без повторений — 6 чисел:
$14,17,41,47,71,74$. С повторениями — 9 чисел, дополнительно
$11,44,77$.

В списке $02,06,20,26,60,62$ исключаются $02$ и $06$; остаются четыре числа.
Они охватывают обе допустимые первые цифры 2 и 6 и обе оставшиеся цифры на
второй позиции.

Условие для $33,35,53,55$: «составить все двузначные числа из цифр 3 и 5;
повторения разрешены». При запрете: $35,53$.

\section{Диагностика ошибок}

\begin{teacherbox}[Не смешивать причины]
  \begin{itemize}
    \item Нет $22$ и $55$ — самовольный запрет повторений.
    \item Есть $038$ — непонимание ведущего нуля.
    \item Верные числа есть, но список неполон — случайный перебор.
    \item Есть число с другой цифрой или длиной — условие не удерживается.
  \end{itemize}
\end{teacherbox}

\section{Контроль и развилка}

Ответы: $11,14,41,44$; $407,470,704,740$.

\begin{resultbox}[Решение]
  \begin{itemize}
    \item \textbf{M06:} списки полны, повторения различены, ведущего нуля нет,
      полнота объяснена группировкой по первой цифре.
    \item \textbf{M05-R:} список случайный, условие о повторениях изменено или
      ноль поставлен первым.
    \item \textbf{M05\(\star\):} ученик может проверять полноту чужого списка
      по структуре, не выписывая всё заново.
  \end{itemize}
\end{resultbox}

После коррекции: цифры 0, 3, 6; составить все двузначные числа без повторений.
Ожидается $30,36,60,63$ с объяснением двух допустимых первых цифр.
